\documentclass[10pt,a4paper]{amsart}
\usepackage{amsmath,amssymb,amsthm,graphicx}
% Retained preamble for the Stacks Project chapter divisors.tex (Divisors).
% The chapter file pinned-original-source.tex carries only \input{preamble} on its first line, so no
% preamble text lives inside the chapter itself.  The Stacks Project uses one single shared preamble for
% every chapter; the snapshot stages it verbatim in sources/stacks-chow/preamble.tex.  The block below is
% copied mechanically, line for line, from that file's lines 14-18 (the Xy-pic setup) and lines 154-217
% (the native theorem declarations together with the mathematical macros).  Nothing was added, removed or
% redefined.  The wrapper class and the amsmath/amssymb/amsthm/graphicx base packages are supplied by the
% assembler, so the class/\externaldocument part of the shared preamble is not repeated here.
% For commutative diagrams we use Xy-pic
\usepackage[all]{xy}

% We use 2cell for 2-commutative diagrams.
\xyoption{2cell}
% Theorem environments.
%
\theoremstyle{plain}
\newtheorem{theorem}[subsection]{Theorem}
\newtheorem{proposition}[subsection]{Proposition}
\newtheorem{lemma}[subsection]{Lemma}

\theoremstyle{definition}
\newtheorem{definition}[subsection]{Definition}
\newtheorem{example}[subsection]{Example}
\newtheorem{exercise}[subsection]{Exercise}
\newtheorem{situation}[subsection]{Situation}

\theoremstyle{remark}
\newtheorem{remark}[subsection]{Remark}
\newtheorem{remarks}[subsection]{Remarks}

\numberwithin{equation}{subsection}

% Macros
%
\def\lim{\mathop{\mathrm{lim}}\nolimits}
\def\colim{\mathop{\mathrm{colim}}\nolimits}
\def\Spec{\mathop{\mathrm{Spec}}}
\def\Hom{\mathop{\mathrm{Hom}}\nolimits}
\def\Ext{\mathop{\mathrm{Ext}}\nolimits}
\def\SheafHom{\mathop{\mathcal{H}\!\mathit{om}}\nolimits}
\def\SheafExt{\mathop{\mathcal{E}\!\mathit{xt}}\nolimits}
\def\Sch{\mathit{Sch}}
\def\Mor{\mathop{\mathrm{Mor}}\nolimits}
\def\Ob{\mathop{\mathrm{Ob}}\nolimits}
\def\Sh{\mathop{\mathit{Sh}}\nolimits}
\def\NL{\mathop{N\!L}\nolimits}
\def\CH{\mathop{\mathrm{CH}}\nolimits}
\def\proetale{{pro\text{-}\acute{e}tale}}
\def\etale{{\acute{e}tale}}
\def\QCoh{\mathit{QCoh}}
\def\Ker{\mathop{\mathrm{Ker}}}
\def\Im{\mathop{\mathrm{Im}}}
\def\Coker{\mathop{\mathrm{Coker}}}
\def\Coim{\mathop{\mathrm{Coim}}}

% Boxtimes
%
\DeclareMathSymbol{\boxtimes}{\mathbin}{AMSa}{"02}

%
% Macros for moduli stacks/spaces
%
\def\QCohstack{\mathcal{QC}\!\mathit{oh}}
\def\Cohstack{\mathcal{C}\!\mathit{oh}}
\def\Spacesstack{\mathcal{S}\!\mathit{paces}}
\def\Quotfunctor{\mathrm{Quot}}
\def\Hilbfunctor{\mathrm{Hilb}}
\def\Curvesstack{\mathcal{C}\!\mathit{urves}}
\def\Polarizedstack{\mathcal{P}\!\mathit{olarized}}
\def\Complexesstack{\mathcal{C}\!\mathit{omplexes}}
% \Pic is the operator that assigns to X its picard group, usage \Pic(X)
% \Picardstack_{X/B} denotes the Picard stack of X over B
% \Picardfunctor_{X/B} denotes the Picard functor of X over B
\def\Pic{\mathop{\mathrm{Pic}}\nolimits}
\def\Picardstack{\mathcal{P}\!\mathit{ic}}
\def\Picardfunctor{\mathrm{Pic}}
\def\Deformationcategory{\mathcal{D}\!\mathit{ef}}

% --- appended 2026-10-12 (Chief): the project own definitions of the reference/slogan/history
% display environments, copied mechanically from the shared Stacks preamble lines 8-12
% (\usepackage{verbatim} plus the three \newenvironment lines, which the chapter preamble keeps
% verbatim from the project preamble). Nothing added, removed or redefined: these environments
% carry no mathematical content and the project itself typesets them as comments. Without these
% three lines a counted statement whose own text uses them fails with Environment slogan undefined.
\usepackage{verbatim}
\newenvironment{reference}{\comment}{\endcomment}
\newenvironment{slogan}{\comment}{\endcomment}
\newenvironment{history}{\comment}{\endcomment}
\usepackage{cleveref}
\crefname{originallemma}{original lemma}{original lemmas}
\Crefname{originallemma}{Original lemma}{Original lemmas}
\makeatletter
\newlabel{lemma-check-reflexive-on-affines}{{0AY0}{source}{}{sourcepointer-1}{}}
\newlabel{lemma-check-reflexive-on-affines@cref}{{[originallemma][0][]0AY0}{[1][0][]source}{}{}{}}
\makeatother
\crefname{originallemma}{original lemma}{original lemmas}
\Crefname{originallemma}{Original lemma}{Original lemmas}
\makeatletter
\newlabel{lemma-reflexive-S2-extend}{{0EBJ}{source}{}{sourcepointer-2}{}}
\newlabel{lemma-reflexive-S2-extend@cref}{{[originallemma][0][]0EBJ}{[1][0][]source}{}{}{}}
\makeatother
\crefname{originallemma}{original lemma}{original lemmas}
\Crefname{originallemma}{Original lemma}{Original lemmas}
\makeatletter
\newlabel{modules-lemma-finite-presentation-stalk-free}{{modules-lemma-finite-presentation-stalk-free}{source}{}{sourcepointer-3}{}}
\newlabel{modules-lemma-finite-presentation-stalk-free@cref}{{[originallemma][0][]modules-lemma-finite-presentation-stalk-free}{[1][0][]source}{}{}{}}
\makeatother
\crefname{originallemma}{original lemma}{original lemmas}
\Crefname{originallemma}{Original lemma}{Original lemmas}
\makeatletter
\newlabel{more-algebra-lemma-dedekind-torsion-free-flat}{{more-algebra-lemma-dedekind-torsion-free-flat}{source}{}{sourcepointer-4}{}}
\newlabel{more-algebra-lemma-dedekind-torsion-free-flat@cref}{{[originallemma][0][]more-algebra-lemma-dedekind-torsion-free-flat}{[1][0][]source}{}{}{}}
\makeatother
\crefname{originallemma}{original lemma}{original lemmas}
\Crefname{originallemma}{Original lemma}{Original lemmas}
\makeatletter
\newlabel{more-algebra-lemma-reflexive-over-normal}{{more-algebra-lemma-reflexive-over-normal}{source}{}{sourcepointer-5}{}}
\newlabel{more-algebra-lemma-reflexive-over-normal@cref}{{[originallemma][0][]more-algebra-lemma-reflexive-over-normal}{[1][0][]source}{}{}{}}
\makeatother
\crefname{originallemma}{original lemma}{original lemmas}
\Crefname{originallemma}{Original lemma}{Original lemmas}
\makeatletter
\newlabel{properties-lemma-codimension-local-ring}{{properties-lemma-codimension-local-ring}{source}{}{sourcepointer-6}{}}
\newlabel{properties-lemma-codimension-local-ring@cref}{{[originallemma][0][]properties-lemma-codimension-local-ring}{[1][0][]source}{}{}{}}
\makeatother
\crefname{originallemma}{original lemma}{original lemmas}
\Crefname{originallemma}{Original lemma}{Original lemmas}
\makeatletter
\newlabel{properties-lemma-criterion-normal}{{properties-lemma-criterion-normal}{source}{}{sourcepointer-7}{}}
\newlabel{properties-lemma-criterion-normal@cref}{{[originallemma][0][]properties-lemma-criterion-normal}{[1][0][]source}{}{}{}}
\makeatother
\title{Reflexivity of a coherent module on an integral normal scheme is torsion freeness plus S2, equivalently extendability from a big open}
\author{The Stacks Project Authors}
\date{stacks-project:divisors.tex:local-snapshot-20261012; GNU FDL 1.2 or later}
\begin{document}
\maketitle
\noindent\textit{Editorial note.} All delivered mathematical text is the author's own, copied from the pinned Stacks Project chapter file pinned-original-source.tex; no statement or proof step was rewritten, reordered or removed.\par\medskip
\noindent\textit{Editorial note.} Following the multi-claim convention for this source, source\_url is the source-level URL https://stacks.math.columbia.edu/ and the source label lemma-reflexive-over-normal is carried in claim\_key and here rather than in a per-claim tag URL.\par\medskip
\noindent\textit{Editorial note.} The preamble is the chapter's shared Stacks Project preamble, retained as sources/stacks-divisors/retained-preamble.tex: a line-for-line excerpt of the shared Stacks preamble staged with this snapshot. No macro was invented or redefined, and this card adds no preamble text of its own.\par\medskip
\noindent\textit{Editorial note.} Statement span 2254-2270 and proof span 2272-2301 come from the halfA pool card operations/pool-cards/stacks-divisors-halfA.json and were re-read in the pinned chapter file before this card was written.\par\medskip
\noindent\textit{Editorial note.} The proof obtains the two directions of (3) with the same cited criterion for normal local rings, applied once for the depth statement outside the big open and once at the end; both citations are preserved and handled as external navigation pointers.\par\medskip
\noindent\textit{Editorial note.} 7 cross-references occur inside the counted spans; they are listed in external\_references as navigation pointers only.\par\medskip
\noindent\textit{Allowed standard prerequisites (editorial).}\begin{itemize}
\item The affine criterion for reflexivity of a coherent module, cited from the same chapter, which reduces reflexivity to the corresponding statement over local rings. Cited at proof lines 2272-2273. Reduces the equivalence of (1) and (2) to the algebraic statement over normal local rings.
\item The extension lemma for coherent modules with property S2 across a closed subset of codimension at least two, cited from the same chapter. Cited at proof lines 2280-2281 and again at 2300-2301. Supplies both directions linking S2 (equivalently sufficient depth outside a big open) with extendability of the restriction from the big open.
\item Local rings of points outside a big open in an integral normal locally Noetherian scheme have depth at least two, and rings of dimension at most one there are fields or discrete valuation rings. Cited at proof lines 2276-2278, 2292-2295 and 2296-2299. Gives the depth and dimension statements that identify the complement of the maximal locally free open as being of codimension at least two.
\item A finite locally free module is reflexive and, over a discrete valuation ring or a field, a torsion free module is free. Used at proof lines 2279-2280 and 2296-2299. Provides the easy inclusion (3) implies (1) and the freeness of the stalks at points of local dimension at most one.
\end{itemize}\medskip
\ifdefined\hypertarget\hypertarget{sourcepointer-1}{}\fi
\noindent\textit{Source reference.} Original lemma 0AY0: focused-seeds/AGP/sources/stacks-snapshot-20261011/tags.txt, source lines 13773--13773. Same-chapter reference used in the retained author text at source line(s) 2273. It is supplied only as a navigation pointer so that the retained text compiles; the counted item delivers the author's own statement and the author's complete proof, and no part of the referenced item is copied in or substituted for the counted argument.\par\medskip
\ifdefined\hypertarget\hypertarget{sourcepointer-2}{}\fi
\noindent\textit{Source reference.} Original lemma 0EBJ: focused-seeds/AGP/sources/stacks-snapshot-20261011/tags.txt, source lines 17924--17924. Same-chapter reference used in the retained author text at source line(s) 2280, 2300, 7349, 7397, 7446, 7506, 7534. It is supplied only as a navigation pointer so that the retained text compiles; the counted item delivers the author's own statement and the author's complete proof, and no part of the referenced item is copied in or substituted for the counted argument.\par\medskip
\ifdefined\hypertarget\hypertarget{sourcepointer-3}{}\fi
\noindent\textit{Source reference.} Original lemma modules-lemma-finite-presentation-stalk-free: sources/stacks-divisors/pinned-original-source.tex, source lines 2288--2288. Same-chapter reference used in the retained author text at source line(s) 2288. It is supplied only as a navigation pointer so that the retained text compiles; the counted item delivers the author's own statement and the author's complete proof, and no part of the referenced item is copied in or substituted for the counted argument.\par\medskip
\ifdefined\hypertarget\hypertarget{sourcepointer-4}{}\fi
\noindent\textit{Source reference.} Original lemma more-algebra-lemma-dedekind-torsion-free-flat: sources/stacks-divisors/pinned-original-source.tex, source lines 1822--1822. Same-chapter reference used in the retained author text at source line(s) 1822, 2293, 2361. It is supplied only as a navigation pointer so that the retained text compiles; the counted item delivers the author's own statement and the author's complete proof, and no part of the referenced item is copied in or substituted for the counted argument.\par\medskip
\ifdefined\hypertarget\hypertarget{sourcepointer-5}{}\fi
\noindent\textit{Source reference.} Original lemma more-algebra-lemma-reflexive-over-normal: sources/stacks-divisors/pinned-original-source.tex, source lines 2275--2275. Same-chapter reference used in the retained author text at source line(s) 2275, 2364. It is supplied only as a navigation pointer so that the retained text compiles; the counted item delivers the author's own statement and the author's complete proof, and no part of the referenced item is copied in or substituted for the counted argument.\par\medskip
\ifdefined\hypertarget\hypertarget{sourcepointer-6}{}\fi
\noindent\textit{Source reference.} Original lemma properties-lemma-codimension-local-ring: sources/stacks-divisors/pinned-original-source.tex, source lines 2295--2295. Same-chapter reference used in the retained author text at source line(s) 2295, 3518, 6583. It is supplied only as a navigation pointer so that the retained text compiles; the counted item delivers the author's own statement and the author's complete proof, and no part of the referenced item is copied in or substituted for the counted argument.\par\medskip
\ifdefined\hypertarget\hypertarget{sourcepointer-7}{}\fi
\noindent\textit{Source reference.} Original lemma properties-lemma-criterion-normal: sources/stacks-divisors/pinned-original-source.tex, source lines 2277--2277. Same-chapter reference used in the retained author text at source line(s) 2277, 2291, 2297, 3050, 7350, 7398, 7412, 7507, 7535. It is supplied only as a navigation pointer so that the retained text compiles; the counted item delivers the author's own statement and the author's complete proof, and no part of the referenced item is copied in or substituted for the counted argument.\par\medskip
\begin{lemma}
\label{lemma-reflexive-over-normal}
Let $X$ be an integral locally Noetherian normal scheme.
Let $\mathcal{F}$ be a coherent $\mathcal{O}_X$-module.
The following are equivalent
\begin{enumerate}
\item $\mathcal{F}$ is reflexive,
\item $\mathcal{F}$ is torsion free and has property $(S_2)$, and
\item there exists an open subscheme $j : U \to X$ such that
\begin{enumerate}
\item every irreducible component of $X \setminus U$
has codimension $\geq 2$ in $X$,
\item $j^*\mathcal{F}$ is finite locally free, and
\item $\mathcal{F} = j_*j^*\mathcal{F}$.
\end{enumerate}
\end{enumerate}
\end{lemma}

\begin{proof}
Using Lemma \ref{lemma-check-reflexive-on-affines}
the equivalence of (1) and (2) follows from
More on Algebra, Lemma \ref{more-algebra-lemma-reflexive-over-normal}.
Let $U \subset X$ be as in (3). By
Properties, Lemma \ref{properties-lemma-criterion-normal}
we see that $\text{depth}(\mathcal{O}_{X, x}) \geq 2$
for $x \not \in U$. Since a finite locally free module is reflexive,
we conclude (3) implies (1) by Lemma \ref{lemma-reflexive-S2-extend}.

\medskip\noindent
Assume (1). Let $U \subset X$ be the maximal open subscheme such
that $j^*\mathcal{F} = \mathcal{F}|_U$
is finite locally free. So (3)(b) holds. Let $x \in X$ be a point.
If $\mathcal{F}_x$ is a free $\mathcal{O}_{X, x}$-module, then
$x \in U$, see
Modules, Lemma \ref{modules-lemma-finite-presentation-stalk-free}.
If $\dim(\mathcal{O}_{X, x}) \leq 1$, then $\mathcal{O}_{X, x}$
is either a field or a discrete valuation ring
(Properties, Lemma \ref{properties-lemma-criterion-normal})
and hence $\mathcal{F}_x$ is free (More on Algebra, Lemma
\ref{more-algebra-lemma-dedekind-torsion-free-flat}).
Thus $x \not \in U \Rightarrow \dim(\mathcal{O}_{X, x}) \geq 2$.
Then Properties, Lemma \ref{properties-lemma-codimension-local-ring}
shows (3)(a) holds. By the already used
Properties, Lemma \ref{properties-lemma-criterion-normal}
we also see that $\text{depth}(\mathcal{O}_{X, x}) \geq 2$
for $x \not \in U$ and hence (3)(c) follows from
Lemma \ref{lemma-reflexive-S2-extend}.
\end{proof}

\end{document}
